\section{Query dan Unifikasi}

Query diawali dengan prompt \texttt{?-} dan berisi tujuan (goal) yang ingin dibuktikan \cite{learn_prolog,swi_prolog,prolog_docs}. Prolog mencari klausa (fakta atau aturan) yang cocok dengan goal dan mencoba memenuhi tujuan tersebut. Jika goal tidak mengandung variabel, jawaban berupa \texttt{true} atau \texttt{false}. Jika goal mengandung variabel, Prolog mengembalikan substitusi (binding) yang membuat goal benar; pengguna dapat menekan \texttt{;} untuk meminta solusi alternatif, atau \texttt{Enter} untuk mengakhiri \cite{learn_prolog,swi_manual}.

\logic{Unifikasi} adalah algoritma pattern-matching yang menemukan substitusi variabel sehingga dua term menjadi identik \cite{learn_prolog,lpn_unification,wiki_prolog}. Misalkan query \texttt{?- suka(X, matematika).} dicocokkan dengan fakta \texttt{suka(ali, matematika).} Unifikasi berhasil dengan substitusi \texttt{X = ali}, karena setelah substitusi kedua term sama. Gambar~\ref{fig:unifikasi-bab08} mengilustrasikan proses ini. Unifikasi memungkinkan Prolog ``mengisi'' variabel secara otomatis saat mencocokkan goal dengan kepala klausa \cite{lpn_unification,learn_prolog}.

\begin{figure}[htbp]
\centering
\begin{tikzpicture}[
  node distance=0.6cm,
  every node/.style={font=\small},
  term/.style={rectangle, draw=black, fill=blue!8, rounded corners},
  sub/.style={rectangle, draw=green!60!black, fill=green!10, rounded corners}
]
  \node[term] (t1) {\texttt{suka(X, matematika)}};
  \node[term, right=2.2cm of t1] (t2) {\texttt{suka(ali, matematika)}};
  \node[sub, below=0.9cm of t1.south, xshift=1.1cm] (s) {\texttt{X = ali}};
  \draw[thick, ->] (t1.south) -- ++(0,-0.35) -| (s.north);
  \draw[thick, ->] (t2.south) -- ++(0,-0.35) -| (s.north);
  \node[font=\footnotesize, below=0.2cm of s] {substitusi};
\end{tikzpicture}
\caption{Unifikasi: goal \texttt{suka(X, matematika)} dan fakta \texttt{suka(ali, matematika)} menjadi identik dengan substitusi \texttt{X = ali} \cite{lpn_unification}}
\label{fig:unifikasi-bab08}
\end{figure}

\logic{Resolution} adalah aturan inferensi utama yang dipakai Prolog \cite{learn_prolog,lpn_unification}. Jika goal cocok dengan fakta, goal terbukti. Jika goal cocok dengan kepala aturan, goal diganti dengan tubuh aturan (konjungsi subgoal). Subgoal dievaluasi satu per satu; jika semuanya terpenuhi, goal asli terbukti. Contoh: untuk \texttt{?- kakek(Ahmad, Ali).}, Prolog mencocokkan dengan \texttt{kakek(X,Z) :- ayah(X,Y), orangtua(Y,Z)} sehingga substitusi \texttt{X=Ahmad, Z=Ali}, lalu membuktikan \texttt{ayah(Ahmad,Y)} dan \texttt{orangtua(Y,Ali)} \cite{learn_prolog,spivey_prolog}.

\logic{Backtracking} terjadi ketika suatu subgoal gagal dibuktikan \cite{learn_prolog,lpn_unification}. Mesin Prolog mundur ke titik pilihan terakhir (choice point), membatalkan unifikasi yang telah dilakukan, dan mencoba klausa atau alternatif lain. Gambar~\ref{fig:resolusi-backtrack-bab08} menyajikan alur skematik: goal $\rightarrow$ unifikasi dengan klausa $\rightarrow$ subgoal; jika subgoal gagal, backtrack ke choice point dan coba klausa lain. Misalnya, jika dua fakta \texttt{suka(ali, matematika)} dan \texttt{suka(budi, matematika)} ada, query \texttt{?- suka(X, matematika).} pertama mengembalikan \texttt{X = ali}; jika pengguna menekan \texttt{;}, Prolog backtrack dan mencoba klausa berikutnya, mengembalikan \texttt{X = budi} \cite{learn_prolog,swi_manual}.

\begin{figure}[htbp]
\centering
\begin{tikzpicture}[
  node distance=0.7cm,
  every node/.style={font=\small},
  box/.style={rectangle, draw=black, fill=blue!10, minimum width=2cm, align=center},
  arrow/.style={thick, ->, >=stealth}
]
  \node[box] (g) {Goal};
  \node[box, below=of g] (u) {Unifikasi dengan\\klausa};
  \node[box, below=of u] (s) {Subgoal};
  \node[box, right=2.2cm of s] (ok) {Terbukti};
  \node[box, left=2.2cm of s] (fail) {Gagal};
  \node[box, below=0.6cm of fail] (bt) {Backtrack\\choice point};
  \draw[arrow] (g) -- (u);
  \draw[arrow] (u) -- (s);
  \draw[arrow] (s) -- (ok);
  \draw[arrow] (s) -- (fail);
  \draw[arrow] (fail) -- (bt);
  \draw[arrow] (bt) |- (u);
\end{tikzpicture}
\caption{Alur eksekusi Prolog: resolution dan backtracking \cite{learn_prolog,lpn_unification}}
\label{fig:resolusi-backtrack-bab08}
\end{figure}

Tabel~\ref{tab:contoh-query-bab08} memberi contoh query dan binding/solusi yang dikembalikan Prolog (dengan program silsilah dari Subsection sebelumnya).

\begin{table}[htbp]
\centering
\small
\begin{tabular}{|>{\raggedright\arraybackslash}p{5.2cm}|>{\raggedright\arraybackslash}p{5.5cm}|}
\hline
\textbf{Query} & \textbf{Contoh jawaban (binding/solusi)} \\
\hline
\texttt{?- suka(ali, matematika).} & \texttt{true} \\
\hline
\texttt{?- suka(X, matematika).} & \texttt{X = ali} ; \texttt{X = budi} (jika ada dua fakta) \\
\hline
\texttt{?- kakek(X, Y).} & \texttt{X = ahmad, Y = ali} ; \texttt{X = ahmad, Y = sari} ; ... \\
\hline
\texttt{?- orangtua(ahmad, X).} & \texttt{X = ali} ; \texttt{X = sari} \\
\hline
\texttt{?- ayah(X, Y), orangtua(Y, Z).} & Pasangan (X,Y,Z) yang memenuhi kedua goal \\
\hline
\end{tabular}
\caption{Contoh query dan hasil (binding) \cite{learn_prolog,swi_manual}}
\label{tab:contoh-query-bab08}
\end{table}

Urutan pencarian Prolog bersifat depth-first dan left-to-right \cite{learn_prolog,spivey_prolog}. Klausa diperiksa sesuai urutan penulisan; literal dalam tubuh aturan dievaluasi dari kiri ke kanan. Urutan ini mempengaruhi efisiensi dan kadang perilaku program (misalnya, infinite loop jika aturan rekursif tidak ditulis dengan hati-hati). Pemahaman urutan evaluasi membantu menulis program yang efisien dan menghindari kegagalan yang tidak diinginkan \cite{learn_prolog,lpn_unification}.

Query majemuk dapat berisi beberapa literal yang dipisahkan koma \cite{learn_prolog,swi_manual}. Contoh \texttt{?- ayah(X,Y), suka(Y, matematika).} menanyakan: ``Siapa ayah dari siapa yang suka matematika?'' Prolog memenuhi keduanya secara berurutan; variabel Y diikat pada literal pertama lalu digunakan di literal kedua. Untuk mencoba semua solusi, tekan \texttt{;} setelah setiap jawaban. Gunakan SWI-Prolog (\url{https://swi-prolog.org/}) atau Learn Prolog Now (\cite{learn_prolog}) untuk praktik interaktif.

\begin{contoh}
\raggedright
Contoh interaksi lengkap: program berisi fakta \texttt{suka(ali, matematika).}, \texttt{suka(budi, matematika).}, dan aturan \texttt{teman(X,Y) :- suka(X,Z), suka(Y,Z), X \textbackslash= Y.} Query dengan variabel dan backtracking (tekan \texttt{;} untuk solusi berikutnya):
\begin{lstlisting}[style=prologstyle, caption={Query dan backtracking}]
?- suka(X, matematika).
X = ali ;
X = budi.

?- teman(ali, X).
X = budi.

?- kakek(Ahmad, Ali).
\end{lstlisting}
Untuk \texttt{?- kakek(Ahmad, Ali).}: Prolog mencocokkan aturan \texttt{kakek}, unifikasi \texttt{X=Ahmad, Z=Ali}, lalu membuktikan \texttt{ayah(Ahmad,Y)} dan \texttt{orangtua(Y,Ali)}. Jika ada fakta yang cocok (misalnya \texttt{orangtua(hasan, ali).} dan \texttt{ayah(ahmad, hasan).}), jawabannya \texttt{true} \cite{learn_prolog,swi_manual}.
\end{contoh}

\begin{konsep}
Query: \texttt{?- goal.} Unifikasi: pattern-matching yang menemukan substitusi agar dua term sama. Resolution: goal cocok fakta $\rightarrow$ terbukti; cocok kepala aturan $\rightarrow$ ganti dengan tubuh aturan. Backtracking: mundur ke choice point saat subgoal gagal, coba alternatif lain. Prolog mencari depth-first, left-to-right.
\end{konsep}
